\documentclass[../main.tex]{subfiles}
\usepackage[utf8]{inputenc}
\usepackage[T1]{fontenc}
\usepackage[indonesian]{babel}

\begin{document}

\chapter{Kriptografi Kunci-Publik: Algoritma RSA}

\begin{learningoutcome}
Setelah mempelajari bab ini, mahasiswa diharapkan mampu:
\begin{itemize}
    \item Menjelaskan konsep dasar kriptografi kunci-publik (asimetris) dan perbedaannya dengan kriptografi kunci-simetri (Sub-CPMK 2.3).
    \item Menguraikan prosedur pembangkitan pasangan kunci publik dan privat pada algoritma RSA (Sub-CPMK 2.3).
    \item Menerapkan rumus enkripsi dan dekripsi RSA untuk mengamankan pesan (Sub-CPMK 2.3).
    \item Menganalisis landasan matematis RSA, termasuk Teorema Euler dan Fungsi Totient Euler (Sub-CPMK 2.3).
    \item Mengevaluasi keamanan RSA berdasarkan kompleksitas pemfaktoran bilangan bulat besar dan ancaman komputasi kuantum (Sub-CPMK 2.3).
    \item Mengidentifikasi risiko serangan \textit{Man-in-the-Middle} (MITM) dalam distribusi kunci publik (Sub-CPMK 2.3).
\end{itemize}
\end{learningoutcome}

\section{Pengantar Kriptografi Kunci-Publik}
Kriptografi kunci-publik, atau kriptografi asimetris, adalah sistem yang menggunakan sepasang kunci yang berbeda: \textbf{kunci publik} (\textit{public key}) dan \textbf{kunci privat} (\textit{private key}). Kunci publik dapat disebarkan secara terbuka, sementara kunci privat harus dijaga kerahasiaannya oleh pemiliknya.

Algoritma RSA (Rivest-Shamir-Adleman), yang diperkenalkan pada tahun 1976, merupakan salah satu implementasi kunci-publik yang paling fundamental dan banyak digunakan di dunia.

\section{Landasan Matematis RSA}
Keamanan RSA didasarkan pada fakta bahwa mengalikan dua bilangan prima besar adalah mudah, tetapi memfaktorkan hasil perkaliannya kembali menjadi dua bilangan prima asalnya adalah sangat sulit secara komputasi.

\subsection{Fungsi Totient Euler}
Fungsi Totient Euler, $\phi(n)$, menentukan jumlah bilangan bulat positif kurang dari $n$ yang relatif prima terhadap $n$. Untuk $n = p \cdot q$ (di mana $p$ dan $q$ adalah bilangan prima), maka:
\[ \phi(n) = (p - 1)(q - 1) \]

\subsection{Teorema Euler}
RSA memanfaatkan Teorema Euler yang menyatakan bahwa jika $\text{PBB}(a, n) = 1$, maka:
\[ a^{\phi(n)} \equiv 1 \pmod{n} \]

\section{Prosedur Algoritma RSA}
Sistem RSA bekerja melalui tiga tahap utama: pembangkitan kunci, enkripsi, dan dekripsi.

\subsection{Pembangkitan Pasangan Kunci}
Untuk membuat pasangan kunci, langkah-langkah berikut dilakukan:
\begin{enumerate}
    \item Pilih dua bilangan prima besar yang berbeda, $p$ dan $q$.
    \item Hitung modulus $n = p \cdot q$.
    \item Hitung nilai $\phi(n) = (p - 1)(q - 1)$.
    \item Pilih bilangan bulat $e$ sebagai kunci publik sedemikian sehingga $1 < e < \phi(n)$ dan $\text{PBB}(e, \phi(n)) = 1$.
    \item Hitung kunci privat $d$ sebagai invers modulo dari $e$ terhadap $\phi(n)$:
    \[ d \equiv e^{-1} \pmod{\phi(n)} \quad \text{atau} \quad e \cdot d \equiv 1 \pmod{\phi(n)} \]
\end{enumerate}
\textbf{Hasil}: Kunci Publik $(e, n)$ dan Kunci Privat $(d, n)$.

\subsection{Proses Enkripsi}
Jika Bob ingin mengirim pesan $m$ kepada Alice, Bob menggunakan kunci publik Alice $(e, n)$:
\begin{enumerate}
    \item Ubah pesan $m$ menjadi integer sedemikian sehingga $0 \le m < n$.
    \item Hitung cipherteks $c$ dengan rumus:
    \[ c = m^e \pmod{n} \]
\end{enumerate}

\subsection{Proses Dekripsi}
Alice menerima cipherteks $c$ dan menggunakan kunci privatnya $(d, n)$ untuk mengembalikan pesan asli:
\[ m = c^d \pmod{n} \]

\section{Analisis Keamanan RSA}
Keamanan RSA terletak pada sulitnya menentukan $d$ jika hanya diketahui $(e, n)$. Untuk mendapatkan $d$, penyerang harus mengetahui $\phi(n)$, yang berarti penyerang harus mampu memfaktorkan $n$ menjadi $p$ dan $q$.

\subsection{Ukuran Kunci dan Komputasi}
Semakin besar $n$, semakin sulit proses pemfaktoran. Saat ini, panjang kunci RSA minimal 1024 bit direkomendasikan, namun 2048 bit atau lebih lebih disarankan untuk keamanan jangka panjang.

\subsection{Kriptografi Hibrida (\textit{Hybrid Cryptography})}
Karena RSA jauh lebih lambat daripada algoritma simetris (seperti AES), dalam praktiknya RSA tidak digunakan untuk mengenkripsi pesan besar. Sebaliknya, digunakan sistem \textbf{Kriptografi Hibrida}:
\begin{enumerate}
    \item Sesi dimulai dengan membangkitkan kunci simetris acak (Kunci Sesi).
    \item Kunci Sesi dienkripsi menggunakan kunci publik penerima (menggunakan RSA).
    \item Pesan utama dienkripsi menggunakan Kunci Sesi (menggunakan AES/DES).
    \item Penerima mendekripsi Kunci Sesi dengan kunci privatnya, lalu menggunakan kunci tersebut untuk mendekripsi pesan.
\end{enumerate}

\subsection{Serangan Man-in-the-Middle (MITM)}
Kelemahan utama distribusi kunci publik adalah kemungkinan adanya penyerang (Carol) yang mencegat komunikasi. Carol dapat mengirimkan kunci publik palsu kepada Alice dan Bob, sehingga Carol dapat mendekripsi semua pesan yang mereka kirimkan satu sama lain tanpa terdeteksi.

\begin{obeactivity}
\textbf{Aktivitas 10.1: Simulasi RSA dengan Bilangan Prima Kecil}
1. Pilih dua bilangan prima kecil, misal $p=3$ dan $q=11$.
2. Hitung $n$ dan $\phi(n)$.
3. Pilih $e$ yang relatif prima terhadap $\phi(n)$.
4. Hitung $d$ menggunakan metode coba-coba atau Algoritma Euclidean.
5. Cobalah mengenkripsi pesan $m=5$ dan dekripsi kembali hasilnya.
6. Diskusikan apa yang terjadi jika $e$ tidak relatif prima terhadap $\phi(n)$.
\end{obeactivity}

\begin{obereflection}
\textbf{Latihan 10.1}
1. Mengapa RSA memerlukan bilangan prima yang sangat besar untuk keamanan praktis?
2. Jika $\phi(n)$ diketahui oleh penyerang, apakah kunci privat $d$ dapat ditemukan dengan mudah? Jelaskan.
3. Apa keuntungan utama dari sistem kriptografi hibrida dibandingkan hanya menggunakan RSA?
\end{obereflection}

\begin{obeassessment}
\textbf{Kuis Bab 10}
1. Jelaskan peran fungsi Totient Euler dalam pembangkitan kunci RSA.
2. Bagaimana prosedur dekripsi RSA mengembalikan pesan asli menggunakan Teorema Euler?
3. Mengapa distribusi kunci publik rentan terhadap serangan \textit{Man-in-the-Middle}?
\end{obeassessment}

\begin{competencychecklist}
\begin{tabularx}{\linewidth}{|X|c|}
\hline
Indikator Kompetensi & Tercapai \\
\hline
Saya dapat membedakan kriptografi simetri dan asimetris & [ ] \\
\hline
Saya dapat menghitung pasangan kunci RSA untuk parameter kecil & [ ] \\
\hline
Saya dapat melakukan proses enkripsi dan dekripsi RSA & [ ] \\
\hline
Saya dapat menjelaskan kaitan pemfaktoran bilangan bulat dengan keamanan RSA & [ ] \\
\hline
Saya dapat menganalisis mekanisme kriptografi hibrida & [ ] \\
\hline
\end{tabularx}
\end{competencychecklist}

\end{document}
